\section{Постановка задачи}
\label{sec:Chapter1} \index{Chapter1}

\subsection{Описание задачи}
\label{subsec:task-description}

Рассматривается задача повышения качества многошаговых рассуждений (multi-hop reasoning) больших языковых моделей посредством использования явной графовой структуры как для оценки, так и для обучения.

Пусть дан вопрос $q$, ответ на который требует последовательного логического вывода через $k$ промежуточных шагов. В графе знаний $\mathcal{G} = (V, E)$, где $V$~--- множество сущностей, $E \subseteq V \times R \times V$~--- множество рёбер с метками отношений из $R$, для вопроса~$q$ выделяется \textit{эталонный подграф} $G^* = (V^*, E^*)$, описывающий корректную цепочку вывода от исходных фактов к финальному ответу.

Модель $\mathcal{M}$, получая на вход вопрос $q$, генерирует текстовый ответ $a$, содержащий цепочку рассуждений. Текстовый ответ преобразуется в \textit{сгенерированный граф} $\hat{G} = (\hat{V}, \hat{E})$ с помощью автоматизированного пайплайна извлечения троек.

Задача состоит в том, чтобы:
\begin{enumerate}
    \setlength{\itemsep}{\smallskipamount}
    \item разработать методы автоматического преобразования текстовых рассуждений в графовое представление и сопоставления с эталонными графами;
    \item определить и формализовать метрики структурного соответствия между $\hat{G}$ и $G^*$;
    \item использовать указанные метрики в качестве функции награды при дообучении модели методами обучения с подкреплением;
    \item экспериментально проверить гипотезу о том, что графовая награда повышает качество рассуждений по сравнению с бинарной оценкой финального ответа.
\end{enumerate}

\textbf{Задачи исследования.}
Для достижения поставленной цели необходимо решить следующие задачи:
\begin{enumerate}
    \setlength{\itemsep}{\smallskipamount}
    \item \textit{Генерация обучающего датасета на основе графа знаний Wikidata.} С помощью внешней LLM сформировать датасет вопросов, требующих многошагового рассуждения, с эталонными подграфами, извлечёнными из Wikidata.
    \item \textit{Определение лучшей базовой модели.} Провести сравнительную оценку (evaluation) нескольких языковых моделей на датасете RuleTaker по графовым метрикам и выбрать наиболее подходящую модель для последующего обучения.
    \item \textit{Двухэтапное обучение выбранной модели.} Провести дообучение модели в два этапа:
    \begin{itemize}
        \item Supervised Fine-Tuning (SFT)~--- адаптация модели к формату структурированных рассуждений на датасетах RuleTaker и Wikidata;
        \item Reinforcement Learning (RL)~--- дообучение с использованием функции награды, учитывающей как корректность итогового ответа, так и структурное соответствие сгенерированного графа рассуждений эталонному подграфу.
    \end{itemize}
    \item \textit{Экспериментальная валидация.} Продемонстрировать, что предложенный подход с графовой наградой повышает качество многошаговых рассуждений по сравнению с бейзлайном, использующим только бинарную оценку финального ответа.
\end{enumerate}

\subsection{Формальные определения}
\label{subsec:formal-definitions}

Введём формальные понятия, используемые на протяжении всей работы. Поскольку в работе задействованы два качественно различных источника данных~--- фактический граф знаний Wikidata и синтетические логические задачи RuleTaker,~--- понятие графа рассуждения определяется для каждого из них отдельно.

\begin{definition}[Граф знаний]
    \label{def:kg}
    Граф знаний~--- ориентированный мультиграф $\mathcal{G} = (V, R, E)$, где $V$~--- конечное множество вершин (сущностей), $R$~--- конечное множество типов отношений, $E \subseteq V \times R \times V$~--- множество направленных рёбер (фактов). Каждый факт представляется тройкой $(s, r, o)$, где $s \in V$~--- субъект, $r \in R$~--- отношение, $o \in V$~--- объект.
\end{definition}

\begin{definition}[Граф знаний-путь (Wikidata)]
    \label{def:kg-path}
    Для вопроса $q$, ответ на который определяется путём длины $k$ в графе знаний, эталонным подграфом называется последовательность связанных троек
    \[
        G^* = \big\{(v_0, r_1, v_1),\, (v_1, r_2, v_2),\, \ldots,\, (v_{k-1}, r_k, v_k)\big\} \subseteq \mathcal{G},
    \]
    где $v_0$~--- стартовая сущность, $v_k$~--- финальная сущность (ответ). Корректный вывод требует последовательного применения всех рёбер пути.
\end{definition}

\begin{definition}[Граф доказательства (RuleTaker)]
    \label{def:proof-graph}
    Для логической задачи RuleTaker графом доказательства называется ориентированный ациклический граф $G^* = (V^*, E^*)$, вершины которого $V^*$~--- высказывания вида <<субъект отношение объект>> (факты, условия правил и промежуточные следствия), а рёбра $E^* \subseteq V^* \times V^*$ задают отношение логического следования <<посылка $\rightarrow$ заключение>>. Вопрос-утверждение является корнем графа (встречается только как заключение), факты контекста~--- листьями (встречаются только как посылки). Поскольку одно заключение может выводиться из нескольких альтернативных наборов посылок, граф доказательства эквивалентно представляется отображением
    \[
        \Pi^* : c \mapsto \{S_1, S_2, \ldots\}, \qquad S_i = \{p^{(i)}_1, p^{(i)}_2, \ldots\},
    \]
    сопоставляющим каждому заключению $c$ список допустимых множеств посылок $S_i$.
\end{definition}

\begin{definition}[Сгенерированный граф рассуждения]
    \label{def:generated-graph}
    Сгенерированный граф рассуждения $\hat{G} = (\hat{V}, \hat{E})$~--- граф, полученный автоматическим извлечением троек (для Wikidata) или пар <<посылка~--- заключение>> (для RuleTaker) из текстового ответа модели $\mathcal{M}$ на вопрос $q$.
\end{definition}

\begin{definition}[Метрики структурного соответствия]
    \label{def:metrics}
    Для пары графов $(\hat{G}, G^*)$ определяются следующие метрики на уровне вершин и рёбер:
    \begin{itemize}
        \item Node Precision: $P_V = \dfrac{|\hat{V} \cap V^*|}{|\hat{V}|}$~--- доля корректных вершин среди сгенерированных;
        \item Node Recall: $R_V = \dfrac{|\hat{V} \cap V^*|}{|V^*|}$~--- доля покрытых эталонных вершин;
        \item Node F1: $F_{1,V} = \dfrac{2 P_V R_V}{P_V + R_V}$~--- гармоническое среднее точности и полноты по вершинам;
        \item Edge Precision: $P_E = \dfrac{|\hat{E} \cap E^*|}{|\hat{E}|}$~--- доля корректных рёбер среди сгенерированных;
        \item Edge Recall: $R_E = \dfrac{|\hat{E} \cap E^*|}{|E^*|}$~--- доля покрытых эталонных рёбер;
        \item Edge F1: $F_{1,E} = \dfrac{2 P_E R_E}{P_E + R_E}$~--- гармоническое среднее точности и полноты по рёбрам;
        \item Graph Edit Similarity: $\text{GES} = 1 - \dfrac{\text{GED}(\hat{G}, G^*)}{\max(|\hat{V}| + |\hat{E}|,\; |V^*| + |E^*|)}$, где $\text{GED}$~--- минимальное число операций вставки, удаления и замены вершин и рёбер, необходимых для преобразования $\hat{G}$ в~$G^*$.
    \end{itemize}
\end{definition}

\subsection{Входные и выходные данные}
\label{subsec:input-output}

\textbf{Входные данные:}
\begin{itemize}
    \item Датасет вопросов, требующих многошагового рассуждения. В работе используются два типа датасетов:
    \begin{itemize}
        \item \textit{RuleTaker}~--- синтетический датасет логических задач с явно заданными правилами вывода и фактами;
        \item \textit{Wikidata-пути}~--- вопросы, ответы на которые определяются путями длины $k$ в графе знаний Wikidata.
    \end{itemize}
    \item Граф знаний (Wikidata) или набор правил и фактов (RuleTaker), из которых извлекаются эталонные подграфы $G^*$.
    \item Предобученная языковая модель (Qwen3-0.6B, Llama~3.x, OLMo-Think, GPT-OSS).
\end{itemize}

\textbf{Выходные данные:}
\begin{itemize}
    \item Фреймворк smart-thinking-llm для проведения экспериментов.
    \item Количественные оценки качества рассуждений моделей по метрикам Node F1, Edge F1 и Graph Edit Similarity.
    \item Дообученная модель (Qwen3-0.6B), прошедшая двухэтапный пайплайн SFT~+~RL с графовой наградой.
    \item Сравнительный анализ эффективности графовой награды по отношению к бинарной оценке финального ответа.
\end{itemize}

\subsection{Требования к решению}
\label{subsec:requirements}

\textbf{Функциональные требования.}

\begin{enumerate}
    \item \textbf{Извлечение графа из текста.} Система должна автоматически преобразовывать текстовые рассуждения модели в набор троек $(s, r, o)$, формирующих граф рассуждения $\hat{G}$.
    \item \textbf{Построение эталонных подграфов.} Для каждого вопроса из датасета должен автоматически строиться эталонный подграф $G^*$ на основе графа знаний или правил вывода.
    \item \textbf{Вычисление метрик.} Система должна вычислять метрики структурного соответствия между $\hat{G}$ и $G^*$.
    \item \textbf{Поддержка нескольких датасетов.} Фреймворк должен поддерживать работу с датасетами RuleTaker и Wikidata.
    \item \textbf{Пайплайн обучения.} Система должна реализовывать двухэтапный пайплайн обучения:
    \begin{itemize}
        \item этап Supervised Fine-Tuning для начальной адаптации модели к формату рассуждений;
        \item этап Reinforcement Learning с использованием графовой функции награды.
    \end{itemize}
    \item \textbf{Графовая функция награды.} Функция награды должна учитывать структурное соответствие сгенерированного графа эталонному, а также корректность финального ответа.
\end{enumerate}

\textbf{Нефункциональные требования.}

\begin{enumerate}
    \item \textbf{Масштабируемость.} Фреймворк должен поддерживать модели размером от 0.6B до 7B параметров.
    \item \textbf{Модульность.} Архитектура системы должна допускать замену отдельных компонентов: модели, датасета, функции награды, алгоритма RL.
    \item \textbf{Документированность.} Код должен быть документирован и снабжён инструкциями по воспроизведению экспериментов.
\end{enumerate}

\subsection{Критерии оценки решения}
\label{subsec:criteria}

Решение считается успешным, если выполнены следующие условия:
\begin{enumerate}
    \item Разработан и функционирует фреймворк smart-thinking-llm, обеспечивающий полный цикл от генерации ответов до вычисления графовых метрик.
    \item Проведён сравнительный анализ не менее четырёх моделей на двух датасетах.
    \item Реализован двухэтапный пайплайн обучения SFT~+~RL с графовой наградой.
    \item Экспериментально продемонстрировано устойчивое повышение метрик качества рассуждений (точности и структурных метрик Node/Edge $F_1$, GES) при использовании графовой награды по сравнению с бейзлайном.
\end{enumerate}